The half-ratio interface itself has a simpler two-step form.  Combining it
with the finite-$q$ source split makes the remaining uniform estimates
explicit.

\begin{lemma}[Two-step cone and folded source ledger; proved reduction]
\label{lem:terminal-modal-static-two-step-ledger}
The local conditions
\eqref{eq:terminal-modal-static-local-half-ratios} are equivalent to
\begin{align}
 d_{m-1}+\frac14d_m&\ge0,
 &d_{m-2}-\frac18d_m&\ge0,
 \label{eq:terminal-modal-static-two-step-frontier}\\
 d_j-\frac14d_{j+2}&\ge0,
 &&0\le j\le m-3.
 \label{eq:terminal-modal-static-two-step-cone}
\end{align}

There is an exact source formula for the rows $0\le j\le m-6$.  Put
$\overline D_n=\bar c_n-\bar c_{n-1}/2$, where $\bar c_n$ is the rescaled
correction from \eqref{eq:terminal-modal-finite-rescaling}, and let
$L_k^\circ$ be the positive correction kernel from
\eqref{eq:terminal-modal-correction-positive-green}, with $L_0^\circ=0$.
For a vector $z$ on the reflected line, write
\[
 \mathcal Q_jz=z(j)-\frac14z(j+2),
 \qquad
 \mathcal R_{m,j,n}z
 =2\mathcal Q_jB_0
 (L_{m-1-n}^\circ-L_{m-n}^\circ)z.
\]
Let $e_a^{\rm ev}$ be the even pair of point masses at $\pm a$ and set
\[
 q_n^{\rm ev}=e_{n-2}^{\rm ev}+2e_{n-1}^{\rm ev}-3e_n^{\rm ev}.
\]
At $n=2$ the coincident seed masses are combined, as in
\eqref{eq:terminal-modal-finite-source-packets}.  If $\mathcal B_{m,j}$ is
the value of $2\mathcal Q_jB_0(\overline D_{m-1}-\overline D_m)$ for the
exact initial pair \eqref{eq:terminal-modal-finite-initial-values} but the
constant inputs $E_n=-3/10$, $F_n=9/40$, then
\begin{equation}
 \frac{d_j-d_{j+2}/4}{q^3\delta^{m-1}}
 =\mathcal B_{m,j}
 +\sum_{n=2}^{m-1}\frac{\varepsilon_n}{4}
   \mathcal R_{m,j,n}q_n^{\rm ev}
 +\sum_{n=2}^{m-1}\mu_n
   \mathcal R_{m,j,n}e_n^{\rm ev}.
 \label{eq:terminal-modal-static-folded-source-ledger}
\end{equation}
Moreover, writing $\mu_n=-q\nu_n$, the exact scalar source obeys
\begin{equation}
 \frac25<\nu_n<\frac{43}{75}
 \qquad(2\le n\le m-1).
 \label{eq:terminal-modal-static-mass-window}
\end{equation}
The following three folded-kernel inequalities would therefore prove
\eqref{eq:terminal-modal-static-two-step-cone} on these rows with strict
margin:
\begin{align}
 \mathcal B_{m,j}&\ge-\frac q{64},
 \label{eq:terminal-modal-static-base-kernel-target}\\
 \sum_{n=2}^{m-1}\frac{\varepsilon_n}{4}
  \mathcal R_{m,j,n}q_n^{\rm ev}&\ge-\frac q{16},
 \label{eq:terminal-modal-static-derivative-kernel-target}\\
 -\sum_{n=2}^{m-1}\nu_n
  \mathcal R_{m,j,n}e_n^{\rm ev}&\ge\frac3{16}.
 \label{eq:terminal-modal-static-mass-kernel-target}
\end{align}
These three uniform kernel bounds and the five direct frontier rows left by
$j>m-6$ remain open.  The formula
\eqref{eq:terminal-modal-static-folded-source-ledger} is an exact reduction,
not evidence that any summand has a pointwise sign.
\end{lemma}

\begin{proof}
By definition,
\[
 E_1=\frac{d_{m-1}+d_m/4}{q^3},\qquad
 F_2-\frac12E_1=\frac{d_{m-2}-d_m/8}{q^3}.
\]
For $r\ge3$, with $j=m-r$, the adjacent terms cancel exactly:
\[
 F_r-\frac12F_{r-1}
 =\frac{d_j-d_{j+2}/4}{q^3}.
\]
This proves the equivalence.

Equation \eqref{eq:terminal-modal-finite-rescaled-target} gives
\[
 \frac{\mathcal T_m}{q^3}
 =\delta^{m-2}(\overline D_{m-1}-\overline D_m).
\]
Apply the interior identity
\eqref{eq:terminal-modal-static-temporal-cone-interior}, followed by
$\mathcal Q_j$.  In the exact source split
\eqref{eq:terminal-modal-finite-source-packets}, a source inserted at time
$n$ reaches $\overline D_N$ through $L_{N-n}^\circ$.  Thus its temporal
difference at $N=m$ is
$(L_{m-1-n}^\circ-L_{m-n}^\circ)$, including $n=m-1$ because
$L_0^\circ=0$.  Superposition of the exact initial/constant-input trajectory,
the derivative packet, and the newest-row mass proves
\eqref{eq:terminal-modal-static-folded-source-ledger}.  The range
$j\le m-6$ is precisely where both $d_j$ and $d_{j+2}$ use the proved
interior/seed smoothing identities.

It remains to prove the lower half of
\eqref{eq:terminal-modal-static-mass-window}; the upper half is
\eqref{eq:terminal-modal-finite-mu-bound}.  From
\eqref{eq:terminal-modal-finite-mu-M},
\[
 \nu_n=\delta^{-(n-1)}\mathfrak m_n,qquad
 \mathfrak m_n=
 \frac{\delta\wp_{n-1}}4R_n
 +\frac{\delta s_n+2\eta}{10\eta},qquad
 R_n=\frac{2\eta-s_n}{q}.
\]
Here $u=\chi^n\le\chi^2$ and
\eqref{eq:terminal-modal-finite-mass-ratio} gives $R_n\ge q$.  Indeed,
after setting $u=\chi^2$, the difference from the required lower bound
factors as
\[
 \frac{2(1-\chi^4)}{1+\chi^4}-(4q-2q^3)
 =\frac{2q(q-1)(q+1)(q^4+5q^2-2)}{q^4+6q^2+1}>0.
\]
Since $s_n\le1$,
\begin{align*}
 \mathfrak m_n-\frac25
 &=R_n\left(\frac{\delta\wp_{n-1}}4-\frac q{10\eta}\right)
   -\frac{qs_n}{10\eta}\\
 &\ge q\left(\frac{\delta\wp_{n-1}}4-\frac{1+q}{10\eta}\right)>0.
\end{align*}
The last inequality uses $\wp_{n-1}>9/5$ from
\cref{lem:terminal-modal-frontier-bounds},
$(1+q)/(10\eta)=1/(5\delta)$, and $9\delta^2>4$.
Finally $\delta^{-(n-1)}>1$, proving $\nu_n>2/5$.

If the three displayed kernel targets hold, the right side of
\eqref{eq:terminal-modal-static-folded-source-ledger} is at least
$-q/64-q/16+3q/16=7q/64>0$.  This proves the stated conditional
implication while leaving the kernel estimates themselves open.
\end{proof}
